\begin{tresc}
Asia, Basia i Kasia bawią się w sklep używając kolorowych żetonów jako pieniędzy. Asia ma w portmonetce 1 żółty żeton, 2 czerwone i 8 niebieskich, Basia ma 1 żółty, 1 czerwony i 10 niebieskich, a Kasia ma po 2 żetony w każdym z trzech kolorów. Jeśli wiemy, że wartość żetonów każdej z dziewczynek jest taka sama, to żółty żeton jest wart tyle co:
\end{tresc}
\begin{wariant}{A} 1 czerwony \end{wariant}
\begin{wariant}{B} 2 czerwone \end{wariant}
\begin{wariant}{C} 3 czerwone \end{wariant}
\begin{wariant}{D} 4 czerwone \end{wariant}
\begin{wariant}{E} 5 czerwonych \end{wariant}
